\section{Related work}
\label{sec:related-work}

Interface design is a separate intervention from the choice of aggregation rule. \citet{Yang2024SWEAgent} study command shape, observations and feedback with a fixed model; their search ablation shows that another available interaction can worsen performance. \citet{Chen2024SAgents} expose delegation and asynchronous communication, but use a hierarchical organization in Minecraft. These studies motivate testing which actions and state a hub exposes. They do not establish that more recruiting or more breakout rooms improve a flat software-engineering team's results. The qualitative follow-up and its measurement boundaries are recorded in \texttt{docs/swarm-harness-intuition-2026-10-03.md}.

An operational follow-up concerns the cost of repeatedly supplying context, rather than the choice of aggregation rule. \citet{Xu2024Compression} retain exact tool and parameter names when selectively compressing tool documentation. This motivates preserving actionable identifiers, but does not establish that the same technique improves hub responses. \citet{Satish2026Compression} show that fewer tokens can accompany more calls and longer completion times, with cache reuse affecting estimated cost. \citet{Liu2026Interaction} find increased information-retrieval calls after compression in a planning environment, but no corresponding retrieval surge in ALFWorld. These results motivate measuring repeated retrieval, cached and uncached input, output tokens, response bytes and task correctness separately. They supply no effect estimate for this hub.

Multi-agent debate provides the immediate setting for this comparison. \citet{Du2023} study debate as a route to better answers, while \citet{Smit2023} examine how debate strategies compare with simpler baselines. Those papers motivate asking whether a conversational protocol earns its additional cost. Our measured question is narrower: on three fixed machine-scored tasks, did this recorded chatroom differ from a vote over fresh attempts?

Arm K's interpreter selector is a form of self-consistency: independently generated answers are normalized and the most frequent is submitted \citep{selfconsistency2022}. Its code selector instead uses execution agreement among candidate programs, following the type of signal studied by \citet{mbrexec2022}. In the observed printf pool, agreement favored the common wrong implementation. That is why the candidate-availability ceiling and the submitted-answer rate must be reported separately. The result describes the selector as used here; it is not a comparison among alternative selectors.



The chatroom's independent openings, challenge, proposal and verification steps make it tempting to attribute corrected answers to a particular interaction. Work on divergent thinking and premature consensus \citep{Liang2023,Yao2025} gives reasons to inspect those interactions. Work on minority answers in multi-agent debate \citep{He2026} likewise frames the wrong-cluster observation as a testable question rather than an established general rule.

Finally, the hidden machine oracle keeps the agents' own consensus separate from correctness. Research on language-model judges \citep{Zheng2023,Song2026} concerns a different measurement route; none of the pass rates here uses such a judge. Our main measurement complication arose earlier, at model serving: the same recorded alias and settings accompanied different thinking-token distributions on different active accounts. The account log and token telemetry therefore belong alongside accuracy and cost in any attempt to reproduce these particular comparisons.
